\section{Conclusion}
\label{sec:conclusion}

This paper addressed the problem of online optimization of bidirectional meeting point
assignment and dynamic routing for many-to-many demand-responsive transit (DRT) with
simulated passenger response via binary logit acceptance.
We formulated the problem as a coupled three-layer model — service offer generation,
binary logit passenger response, and dynamic vehicle dispatch — and solved it
with an ex-post MILP routing diagnostic for small instances and a
rolling horizon Adaptive Large Neighborhood Search (ALNS) heuristic for large-scale
deployment.
Numerical experiments on synthetic scenarios calibrated to Chinese suburban conditions
demonstrated that the full bidirectional model achieves
significantly lower per-trip vehicle utilization than door-to-door service at equal
served share (endogenous matched-coverage comparison).
Under the unconstrained comparison (FullModel 18.3\% vs.\ DoorToDoor 61.0\%), FullModel achieves
15.1\,vkm/trip versus 21.3\,vkm/trip for DoorToDoor (29.1\% improvement).
This result confirms that jointly optimizing both pickup meeting points ($M_r^P$) and
dropoff meeting points ($M_r^D$) — while explicitly accounting for passenger heterogeneity
in mode choice — yields substantial operational efficiency gains over conventional
door-to-door and single-sided meeting point approaches.

The paper made four contributions to the DRT literature.
First, we proposed the first formulation combining bidirectional meeting
point sets, binary logit passenger choice with heterogeneous types, and
rolling horizon re-optimisation for many-to-many DRT — a coupled three-layer model extending prior work
that considered only single-sided meeting points \citep{cortenbach2024} or omitted
passenger choice \citep{wu2025}.
Second, we developed an ex-post MILP routing diagnostic and a rolling horizon ALNS heuristic
specifically designed for the bidirectional dial-a-ride problem with endogenous demand.
Third, we empirically demonstrated efficiency gains in per-trip vehicle utilization
(endogenous matched-coverage comparison) and 29.1\% gain (unconstrained comparison) and conducted
ablation experiments that isolated the contributions of bidirectional meeting points,
rolling horizon scheduling, and passenger choice modeling.
Fourth, we performed an equity analysis across three passenger types, revealing a Gini
coefficient of 0.1216 in acceptance rates: time-sensitive passengers are systematically
disadvantaged (14.4\% acceptance) relative to price-sensitive (25.4\%) and walk-sensitive
(20.2\%) passengers, a finding with direct implications for service equity in Chinese
low-density urban areas.

The policy implications of these findings are organized into five actionable
recommendations for Chinese city DRT operators.
In our simulation setting, bidirectional meeting point deployment yields meaningful
efficiency gains at walking radius $\rho \geq 1000$~m, and fleet ratios below
15 vehicles per 100 peak-hour requests lead to service degradation due to vehicle
scarcity. These scenario-specific findings, together with the three-tier deployment
framework, provide simulation-based guidance for DRT operators and urban
planners designing or scaling DRT services in low-density Chinese urban areas.

Several limitations of this study should be acknowledged.
The binary logit utility parameters ($\beta$ coefficients for fare, waiting time, and walking
distance) were calibrated from simulation rather than from real passenger survey data;
future work should validate these parameters against stated or revealed preference surveys
from Chinese DRT users.
The numerical experiments rely on synthetic and semi-realistic scenarios; validation
against real GPS trajectories and smart card data from operational Chinese DRT systems
remains an important direction for future research.
Additionally, the set of candidate meeting point locations $M$ is treated as a given
input; the joint optimization of meeting point placement and vehicle routing is outside
the scope of this paper and represents a natural extension.

Building on these limitations, we identify four directions for future research.
First, real-data validation using GPS and smart card data from Chinese city DRT pilots
would strengthen the empirical grounding of the model.
Second, anticipatory and lookahead extensions to the rolling horizon framework — such as
scenario-based stochastic optimization \citep{bent2004} — could improve performance
under high demand uncertainty.
Third, incorporating time-dependent travel times would make the model more realistic for
peak-hour DRT operations in congested urban networks.
Fourth, the joint optimization of meeting point placement and dynamic routing represents
a promising research frontier that would unify the location and scheduling decisions
currently treated separately in the literature.
